% ************************************************************
% DISCUSSION
% ************************************************************
\section{Discussion}

The first-pass index classified all analytic strata as high priority, reflecting convergence across WASH service stress, environmental-health hazards, health-system constraints, population vulnerability, and operational constraints. The stability of classifications across weighting specifications suggests that the result is not dependent on a single scoring assumption. The most defensible interpretation is that the available public data identify clear areas for follow-up validation and recovery-prioritization analysis.

This framework is aligned with health-equity and health-systems research because it focuses on variation in exposure, service access, and recovery need. Its practical value lies in converting fragmented public reporting into an auditable prioritization structure that can be updated as more precise geospatial, facility-functionality, and WASH dashboard data become available.
% ************************************************************
% END DISCUSSION
% ************************************************************
